\documentclass[12pt, a4paper]{article}

%%%%%%%%%%%%%%%%%%%%%%%%%%%%%%%%%%%%%%%%%%%%%%%%%%%%%%%%%%%%%%%%%
% Load style file
%%%%%%%%%%%%%%%%%%%%%%%%%%%%%%%%%%%%%%%%%%%%%%%%%%%%%%%%%%%%%%%%%
\input{general.tex}

%%%%%%%%%%%%%%%%%%%%%%%%%%%%%%%%%%%%%%%%%%%%%%%%%%%%%%%%%%%%%%%%%
% Begin document
%%%%%%%%%%%%%%%%%%%%%%%%%%%%%%%%%%%%%%%%%%%%%%%%%%%%%%%%%%%%%%%%%
\begin{document}

\thispagestyle{empty}

%%%%%%%%%%%%%%%%%%%%%%%%%%%%%%%%%%%%%%%%%%%%%%%%%%%%%%%%%%%%%%%%%
% Define the title page
%%%%%%%%%%%%%%%%%%%%%%%%%%%%%%%%%%%%%%%%%%%%%%%%%%%%%%%%%%%%%%%%%

\begin{minipage}{0.4\textwidth}
\begin{flushleft}
	\includegraphics[width=6cm]{example-image} % Add logo of the University
\end{flushleft}
\end{minipage}
\hfill
\begin{minipage}{0.4\textwidth}
\begin{flushright}
\small{Department name, Chair name} % Add department and chair info
\end{flushright}
\end{minipage}

\vspace*{1cm}

\begin{center}
{\Large \textnormal{COURSE TITLE}} \\
\vspace{.25cm}
SEMESTER 
\end{center}

\begin{center}
\rule{6.6in}{0.4pt}
\begin{minipage}[t]{\textwidth}
\begin{tabular}{llcccll}
\textbf{Instructor:} & Instructor name & & &  & \textbf{Time:} & Time of the seminar   \\ 
\textbf{E-mail:} &  \href{mailto:instructor.email@uni.edu}{instructor.email@uni.edu} & & & & \textbf{Room:} & Room Number
\end{tabular}
\end{minipage}
\rule{6.6in}{0.4pt}
\end{center}
\vspace{.5cm}
\setlength{\unitlength}{1in}
\renewcommand{\arraystretch}{2}
\vskip.05in
\noindent\textbf{Course description:}  
This is an exemplary syllabus that can be easily adjusted toward the specific needs of the instructor. It contains some text chunks and is mostly populated with blind text.

\vskip.15in
\noindent\textbf{Requirements:} Download and install \href{https://www.r-project.org}{\textsf{R}} and \href{https://www.rstudio.com}{\textsf{R} Studio}. Throughout the class, we will mainly rely on RStudio Cloud (a browser based version of R). To follow the course, please \href{https://rstudio.cloud/}{sign up here (free version)}.

\vskip.15in
\noindent\textbf{Student hours:} Upon request. Please contact me via e-mail (including a short description of your problem or question). 


%%%%%%%%%%%%%%%%%%%%%%%%%%%%%%%%%%%%%%%%%%%%%%%%%%%%%%%%%%%%%%%%%

\newpage

%%%%%%%%%%%%%%%%%%%%%%%%%%%%%%%%%%%%%%%%%%%%%%%%%%%%%%%%%%%%%%%%%
% Define rules and regulations
%%%%%%%%%%%%%%%%%%%%%%%%%%%%%%%%%%%%%%%%%%%%%%%%%%%%%%%%%%%%%%%%%

\vskip.15in
\noindent\textbf{Important Deadlines:}
\begin{center} \begin{minipage}{6.5in}
\begin{flushleft}
1) Complete ``Intro to \textsf{R}'' (online course)\dotfill 
~04.10.2020 \\
2) Research question \dotfill ~04.10.2020\\
3) Homework \dotfill ~Respective week\\
4) Presentation of course projects  \dotfill ~02.12.2020 (in class)\\
6) Term paper    \dotfill ~ 10.01.2021
\end{flushleft}
\end{minipage}
\end{center}
\vskip.15in

Submissions (1), (2), (3), (5) and (6) must be handed in electronically at 23:59 (the latest) on the respective day.\\
Giving a brief presentation of the course project (4) is expected in-class.
\vskip.15in

\noindent\textbf{Late submission policy:}  \blindtext

\noindent\textbf{Plagiarism:} \blindtext

\noindent\textbf{Grading policy:}   \blindtext

\noindent\textbf{Class attendance and participation:}  The learning process of this class is based on in-class discussion and participation. Attendance and careful preparation of the course material is therefore highly recommended. This includes coming to class on time.


\noindent\textbf{Laptop and phone policy:} In order to ensure an active participation and to keep your attention on the important things (our class), please avoid distracting yourself through (unnecessary) electronic devices or applications. For further insights on the consequences of multitasking, I recommend the study by \textcite{bellur2015make} (\url{https://bit.ly/2GnyTf2}). They found that in-class multitasking leads to significantly lower performance.
\vskip.15in
\noindent\textbf{Interesting links:} There is a number of great websites and podcasts that provide additional interesting information. See for instance \href{https://www.r-bloggers.com}{\textsf{R}-bloggers}, \href{https://www.mzes.uni-mannheim.de/socialsciencedatalab/}{Methods Bites}, \href{https://www.r-graph-gallery.com/portfolio/ggplot2-package/}{\textsf{R} Graph Gallery}, \href{https://towardsdatascience.com}{Towards Data Science},  \href{https://fivethirtyeight.com}{FiveThirtyEight}, and \href{https://rweekly.org}{R Weekly}.

%%%%%%%%%%%%%%%%%%%%%%%%%%%%%%%%%%%%%%%%%%%%%%%%%%%%%%%%%%%%%%%%%
% Course outline
%%%%%%%%%%%%%%%%%%%%%%%%%%%%%%%%%%%%%%%%%%%%%%%%%%%%%%%%%%%%%%%%%
\newpage

\noindent\textbf{Course outline:}  
\begin{itemize}
\item[] The optional readings consist of book chapters. To make them as accessible as possible for interested readers, I will upload some of the readings online.
\item[] I ask you to submit your assignments (\faExclamation) within the deadline and to prepare the \textbf{one} required reading (\faBook) carefully. CheatSheets (\faLightbulbO) and additional input (\faEdit) are optional. All the resources should help you if you need further input or want to expand your knowledge. 
\item[]

%%%%%%%%%%%%%%%%%%%%%%%%%%%%%%%%%%%%%%%%%%%%%%%%%%%%%%%%%%%%%%%%%
% Preparation
%%%%%%%%%%%%%%%%%%%%%%%%%%%%%%%%%%%%%%%%%%%%%%%%%%%%%%%%%%%%%%%%%

\vspace*{0.04in}
\noindent\rule{16cm}{0.7pt}
\vspace*{-0.45in}
\item[] \textbf{Week 0: Preparation}
\item[] \textit{Before starting with the course, I want you to already get some first insights into \textsf{R} so that we can depart from a similar level.}
\item[] \faExclamation \space 
\space Online course: Intro to \textsf{R}
\begin{itemize}   
\item[$\square$] Complete the class ``Intro to \textsf{R}'' (preferably \textbf{before} we meet but\textbf{ no later than DATE}).
\end{itemize}

%%%%%%%%%%%%%%%%%%%%%%%%%%%%%%%%%%%%%%%%%%%%%%%%%%%%%%%%%%%%%%%%%
% Session 1
%%%%%%%%%%%%%%%%%%%%%%%%%%%%%%%%%%%%%%%%%%%%%%%%%%%%%%%%%%%%%%%%%
\vspace*{0.04in}
\noindent\rule{16cm}{0.7pt}
\vspace*{-0.45in}
% Add a short description
\item[] \textbf{Session 1: Introduction to the seminar and intro to R I (DATE)} 
\item[]\textit{This session will introduce you to the seminar in general, give a brief outline of the course, and address organizational issues. We will then have a short introduction to \textsf{R} and \textsf{R} Studio and meet \textsf{R}Markdown -- an efficient way of typesetting text and code in \textsf{R}.\\
We will use the paper by \textcite{hultman2014beyond} as a working example that guides us through this course -- please read it carefully!}
% Add required readings
\item[] \faBook \space Required readings
\begin{itemize} 
\item[$\square$] \fullcite{hultman2014beyond}
\end{itemize}
{\color{darkgray}\begin{footnotesize}
% Add optional reading
\item[] Optional readings
\begin{itemize}
\item[$\ast$] \fullcite{healy2018data}, Chapter 2.1-2.4
\item[$\ast$] \fullcite{imai2018quantitative}, Chapter 1.3
\item[$\ast$] \fullcite{wickham2016r}, Chapter 2, 6, 9, 21, 23, 24
\end{itemize}
\end{footnotesize}}
% Add additional input (YouTube videos, ...)
\item[] \faEdit \space Additional input 
\begin{itemize} 
\item[$\square$] \href{url}{Link}
\end{itemize}
% Add CheatSheets
\item[] \faLightbulbO \space CheatSheet
\begin{itemize} 
\item[$\square$] \href{https://github.com/rstudio/cheatsheets/raw/master/rstudio-ide.pdf}{\textsf{R} Studio}
\item[$\square$] \href{https://github.com/rstudio/cheatsheets/raw/master/rmarkdown-2.0.pdf}{\textsf{R} Markdown}
\item[$\square$] \href{http://442r58kc8ke1y38f62ssb208-wpengine.netdna-ssl.com/wp-content/uploads/2015/03/rmarkdown-reference.pdf}{\textsf{R} Markdown Reference Guide}
\item[$\square$] \href{https://bit.ly/31qRe28}{Tidyverse for Beginners}
\end{itemize}
% Remind students of submissions
\item[] \faExclamation \space 
\space Research question

%%%%%%%%%%%%%%%%%%%%%%%%%%%%%%%%%%%%%%%%%%%%%%%%%%%%%%%%%%%%%%%%%
% Session 2
%%%%%%%%%%%%%%%%%%%%%%%%%%%%%%%%%%%%%%%%%%%%%%%%%%%%%%%%%%%%%%%%%
\vspace{0.5cm}
\item[] \textbf{Session 2: Intro to \textsf{R} II (DATE):}
% Add a short description
\item[] \textit{Now that we have met \textsf{R}, we dedicate this session to the learn more of \textsf{R}'s logic and basics --  objects, vectors, and data frames.}
% Add additional input (YouTube videos, ...)
\item[] \faEdit \space Additional input
\begin{itemize} 
\item[$\square$] \href{url}{Link}
\end{itemize}
% Add CheatSheets
\item[] \faLightbulbO \space CheatSheet
\begin{itemize} 
\item[$\square$] \href{https://ugoproto.github.io/ugo_r_doc/R-Cheat-Sheet-Brief-Introduction-to-Language-Elements-and-Control-Structures.pdf}{Brief Introduction to Language Elements and Control Structures}
\item[$\square$] \href{https://www.rstudio.com/wp-content/uploads/2016/05/base-r.pdf}{Base \textsf{R}}
\end{itemize}

%%%%%%%%%%%%%%%%%%%%%%%%%%%%%%%%%%%%%%%%%%%%%%%%%%%%%%%%%%%%%%%%%
% Session 3
%%%%%%%%%%%%%%%%%%%%%%%%%%%%%%%%%%%%%%%%%%%%%%%%%%%%%%%%%%%%%%%%%
\item[] \textbf{Session 3: Intro to \textsf{R} III (DATE):}
% Add a short description
\item[] \textit{In this session, we will learn how to import different data formats in \textsf{R} and have a first (descriptive) look at the data.}
% Add additional input (YouTube videos, ...)
\item[] \faEdit \space Additional input
\begin{itemize} 
\item[$\square$] \href{url}{Link}
\end{itemize}
% Add CheatSheets
\item[] \faLightbulbO \space CheatSheet
\begin{itemize} 
\item[$\square$] \href{https://ugoproto.github.io/ugo_r_doc/data-import.pdf}{Data Import}
\end{itemize}
\item[] \faExclamation \space Homework
\begin{itemize} 
\item[$\square$] Homework 1 is due today.
\end{itemize}

%%%%%%%%%%%%%%%%%%%%%%%%%%%%%%%%%%%%%%%%%%%%%%%%%%%%%%%%%%%%%%%%%
% Session 4
%%%%%%%%%%%%%%%%%%%%%%%%%%%%%%%%%%%%%%%%%%%%%%%%%%%%%%%%%%%%%%%%%
\vspace*{0.04in}
\noindent\rule{16cm}{0.7pt}
\vspace*{-0.45in}
\item[] \textbf{Session 4: Data wrangling and merging I (DATE)} 
% Add a short description
\item[]\textit{Now that we successfully loaded all data and had some first looks at the descriptive statistics, we will learn how to clean (more or less) messy data in \textsf{R}.}
% Add additional input (YouTube videos, ...)
\item[] \faEdit \space Additional input
\begin{itemize} 
\item[$\square$] \href{url}{Link}
\end{itemize}
% Add CheatSheets
\item[] \faLightbulbO \space CheatSheet
\begin{itemize} 
\item[$\square$] \href{https://www.rstudio.com/wp-content/uploads/2015/02/data-wrangling-cheatsheet.pdf}{Data Wrangling With \texttt{dplyr} and \texttt{tidyr}}
\item[$\square$] \href{https://github.com/cosimameyer/overviewR/blob/master/man/figures/CheatSheet_overviewR.pdf}{\texttt{overviewR}}
\end{itemize}

%%%%%%%%%%%%%%%%%%%%%%%%%%%%%%%%%%%%%%%%%%%%%%%%%%%%%%%%%%%%%%%%%
% Session 5
%%%%%%%%%%%%%%%%%%%%%%%%%%%%%%%%%%%%%%%%%%%%%%%%%%%%%%%%%%%%%%%%%
\item[] \textbf{Session 5: Data wrangling and merging II (DATE)} 
% Add a short description
\item[]\textit{We then proceed and learn how to manipulate data in \textsf{R} to adjust it to our (format) needs before we learn how to merge the data in the next session.}
% Add additional input (YouTube videos, ...)
\item[] \faEdit \space Additional input
\begin{itemize} 
\item[$\square$] \href{url}{Link}
\end{itemize}
% Add CheatSheets
\item[] \faLightbulbO \space CheatSheet
\begin{itemize} 
\item[$\square$] \href{https://www.rstudio.com/wp-content/uploads/2015/02/data-wrangling-cheatsheet.pdf}{Data Wrangling With \texttt{dplyr} and \texttt{tidyr}}
\end{itemize}

%%%%%%%%%%%%%%%%%%%%%%%%%%%%%%%%%%%%%%%%%%%%%%%%%%%%%%%%%%%%%%%%%
% Session 6
%%%%%%%%%%%%%%%%%%%%%%%%%%%%%%%%%%%%%%%%%%%%%%%%%%%%%%%%%%%%%%%%%
\item[] \textbf{Session 6: Data wrangling and merging III (DATE)} 
% Add a short description
\item[]\textit{Merging data can easily become one of the most complex parts of data processing. This session serves as a buffer session..}
% Add additional input (YouTube videos, ...)
\item[] \faEdit \space Additional input
\begin{itemize} 
\item[$\square$] \href{url}{Link}
\end{itemize}
% Add CheatSheets
\item[] \faLightbulbO \space CheatSheet
\begin{itemize} 
\item[$\square$] \href{https://www.rstudio.com/wp-content/uploads/2015/02/data-wrangling-cheatsheet.pdf}{Data Wrangling With \texttt{dplyr} and \texttt{tidyr}}
\end{itemize}
% Remind students of submissions
\item[] \faExclamation \space Homework
\begin{itemize} 
\item[$\square$] Homework 2 is due today.
\end{itemize}

%%%%%%%%%%%%%%%%%%%%%%%%%%%%%%%%%%%%%%%%%%%%%%%%%%%%%%%%%%%%%%%%%
% Session 7
%%%%%%%%%%%%%%%%%%%%%%%%%%%%%%%%%%%%%%%%%%%%%%%%%%%%%%%%%%%%%%%%%
\vspace*{0.04in}
\noindent\rule{16cm}{0.7pt}
\vspace*{-0.45in}

\item[] \textbf{Session 7: Data visualization I (DATE)} 
% Add a short description
\item[]\textit{In the third part of our seminar, we delve into the magic world of \texttt{gpplot2} -- a great way of plotting your results in \textsf{R}. In the first session, we will learn the basics of \texttt{ggplot2}.}
{\color{darkgray}\begin{footnotesize}
% Add optional reading
\item[] Optional readings
\begin{itemize} 
\item[$\ast$] \fullcite{healy2018data}, Chapter 3.1-3.7
\end{itemize}
\end{footnotesize}}
% Add additional input (YouTube videos, ...)
\item[] \faEdit \space Additional input
\begin{itemize} 
\item[$\square$] \href{url}{Link}
\end{itemize}
% Add CheatSheets
\item[] \faLightbulbO \space CheatSheet
\begin{itemize} 
\item[$\square$] \href{https://www.rstudio.com/wp-content/uploads/2015/03/ggplot2-cheatsheet.pdf}{\texttt{ggplot2}}
\end{itemize}

%%%%%%%%%%%%%%%%%%%%%%%%%%%%%%%%%%%%%%%%%%%%%%%%%%%%%%%%%%%%%%%%%
% Session 8
%%%%%%%%%%%%%%%%%%%%%%%%%%%%%%%%%%%%%%%%%%%%%%%%%%%%%%%%%%%%%%%%%
\vspace{-0.1cm}
\item[] \textbf{Session 8: Data visualization II (DATE)} 
% Add a short description
\item[]\textit{In the second session, we will then produce more advanced figures with \texttt{ggplot2}. This \href{https://towardsdatascience.com/tableau-esque-drag-and-drop-gui-visualization-in-r-901ee9f2fe3f}{Medium blogpost} gives you an additional idea how to produce Tableau-like graphs in \textsf{R}.} 
{\color{darkgray}\begin{footnotesize}
% Add optional reading
\item[] Optional readings
\begin{itemize}
\item[$\square$] \fullcite{healy2018data}, Chapter 8.1-8.4
\item[$\ast$] \fullcite{healy2018data}, Chapter 7
\end{itemize}
\end{footnotesize}}
% Add additional input (YouTube videos, ...)
\item[] \faEdit \space Additional input
\begin{itemize} 
\item[$\square$] \href{url}{Link}
\end{itemize}
% Add CheatSheets
\item[] \faLightbulbO \space CheatSheet
\begin{itemize} 
\item[$\square$] \href{https://images.plot.ly/plotly-documentation/images/r_cheat_sheet.pdf}{\texttt{plotly}}
\end{itemize}

%%%%%%%%%%%%%%%%%%%%%%%%%%%%%%%%%%%%%%%%%%%%%%%%%%%%%%%%%%%%%%%%%
% Session 9
%%%%%%%%%%%%%%%%%%%%%%%%%%%%%%%%%%%%%%%%%%%%%%%%%%%%%%%%%%%%%%%%%
\vspace{0.5cm}
% Add a short description
\item[] \textbf{Session 9: Data visualization III and hackathon (DATE)} 
\item[]\textit{We will dedicate this session with for remaining parts of data visualization and you consult me individually for help and address open questions as well as structure your final projects.}
% Remind students of submissions
\item[] \faExclamation \space Homework
\begin{itemize} 
\item[$\square$] Homework 3 is due today.
\end{itemize}

%%%%%%%%%%%%%%%%%%%%%%%%%%%%%%%%%%%%%%%%%%%%%%%%%%%%%%%%%%%%%%%%%
% Session 10
%%%%%%%%%%%%%%%%%%%%%%%%%%%%%%%%%%%%%%%%%%%%%%%%%%%%%%%%%%%%%%%%%
\noindent\rule{16cm}{0.7pt}
\vspace*{-0.45in}
\item[] \textbf{Session 10: Mini conference  (DATE)} 
% Add a short description
\item[]\textit{During this session, we will set up a mini conference setting. You will have the chance to present your course projects and receive constructive feedback from your peers that should help you when preparing your final paper.}

%%%%%%%%%%%%%%%%%%%%%%%%%%%%%%%%%%%%%%%%%%%%%%%%%%%%%%%%%%%%%%%%%
% Session 11
%%%%%%%%%%%%%%%%%%%%%%%%%%%%%%%%%%%%%%%%%%%%%%%%%%%%%%%%%%%%%%%%%
\vspace*{0.04in}
\noindent\rule{16cm}{0.7pt}
\vspace*{-0.45in}
\item[] \textbf{Session 11: Outlook and wrap-up: What else can we do in \textsf{R}? (DATE)} 
% Add a short description
\item[]\textit{To wrap up the seminar, I will give you a some insights what other cool  things that you can do in and with \textsf{R}.}
\item[] \textit{We will also have a final discussion and deal with potential challenges of your term papers. If you have any further questions or would like to discuss something particular, please feel free to contact me beforehand.}
\end{itemize}

%%%%%%%%%%%%%%%%%%%%%%%%%%%%%%%%%%%%%%%%%%%%%%%%%%%%%%%%%%%%%%%%%
% Assessment criteria catalogue
%%%%%%%%%%%%%%%%%%%%%%%%%%%%%%%%%%%%%%%%%%%%%%%%%%%%%%%%%%%%%%%%%
\newpage
\appendix

\subsection*{Assessment criteria}

\subsubsection*{Coursework}
\subsubsection*{1. Complete ``Intro to \textsf{R}''}
\begin{itemize}
\item[] Before we delve into \textsf{R}, I want you to familiarize yourself with first basics in \textsf{R}. 
\end{itemize}

\subsubsection*{2. Research question}
\begin{itemize}
\item[] To match your interests for the course projects, I ask you to submit your research questions. I will then use your research questions to group you according to your interests. Collaboration is key and group works are meant to build these synergistic learning environments where you can learn and teach from and with others. While everyone has to submit her/his individual final paper (project), you should tackle the basic (data) problems in your group together and then use your output to answer your individual research question.\\
I will also provide you with feedback on the feasibility and plausibility of your research question.
 \item[] \textbf{Deadline:} \textbf{\textbf{DATE} \textbf{TIME}}
 \item[] \textbf{Submission:} Online
  \item[] \textbf{Should be included:}
  \begin{itemize}
      \item[-] Analytical research question 
      \item[-] (You may also add a few explanatory sentences.) 
  \end{itemize}
      \end{itemize}
      
\subsubsection*{3. Homework}
\begin{itemize}
\item[] I will assign homework. The homework should help you to strengthen your previously learned knowledge. To pass this coursework, \textbf{you need to pass 2 out of 3} homework.
\item[] \textbf{Deadline:} Respective week
\item[] \textbf{Submission:} Online
\end{itemize}      

\subsubsection*{4. Mini conference -- Presentation of course projects}
\begin{itemize}
\item[] \noindent The presentations give you the chance to present your projects and to receive constructive feedback from your peers.\\
The data presentations should not be longer than 15 minutes (\textbf{maximum}) with slides (PowerPoint, markdown (html, PDF), \LaTeX~(PDF)).
\item[] \textbf{Deadline:} \textbf{DATE} (in class)
\item[] \textbf{Formal requirements:} 15 minutes presentation with slides
\item[] \textbf{Submission:} Presentation (in-class)
\item[] \textbf{Should be included:}
\begin{itemize}
\item[-] What is your group's over-arching research question/thematic focus? 
\item[-] What are your research questions?
\item[-] Which data did you use?
\item[-] What is the data about?
\item[-] What is the unit of analysis and the time frame?
\item[-] How was the data collected?
\item[-] What does the data tell you about your research question(s)? (Come up with some descriptive statistics, some visualization, ...)
\item[-] Where can we access the data? (Reference, format) 
\item[-] Brief concluding discussion
\item[-] Potential problems
\end{itemize}
\end{itemize}

\newpage
\subsubsection*{Final examination}
\subsubsection*{Term paper}
The deadline for the term paper (in English) is \textbf{DATE} (\textbf{TIME}; upload as PDF). \\
For the purpose of consistency, please use a \textit{coherent citation style} (see for example the \href{http://bit.ly/1J6qOrD}{APSA citation guidelines}).\\
\vspace*{0.1in}\\
\textbf{General requirements}
\begin{enumerate}
\item Formal requirements
\begin{itemize}
\item[-] Cover page: university, chair, semester, course type and title, name of the instructor, title of the term paper, date; name, contact information, program and semester of the student; indicate word count
\item[-] Table of contents: complete structure of the final paper including page numbers and informative headlines
\item[-] Length: 5 pages in total (max. 2,000 words, will be checked); 2,000 words include only the content (\textit{no} cover page, table of content, references, or appendices) %; 25 pages (max. 6.000 words) if co-authored
\item[-] Signed statutory declaration (both in English and in German, as uploaded on ILIAS)
\item[-] Layout: 1.5 line spacing, font size 11, Times New Roman, pagination of the text
\item[-] Orthography and grammar
\item[-] Formatted as a PDF or html
\item[-] \textbf{Submit your data files and \textsf{R}markdown file as well}
\end{itemize}
\item Scientific standards
\begin{itemize}
\item[-] Appropriate use and formal correctness of references, e.g. for example according to APSA citation style
\item[-] Independent research
\item[-] List of references: coherent citation style, e.g. for example according to APSA citation style 
\end{itemize}
\end{enumerate}

\pagebreak
\textbf{Content}

\begin{enumerate}
\item Introduction
\begin{itemize}
\item[-] Relevance
\item[-] Analytical research question
\item[-] Overview of the paper's structure
\end{itemize}
\item Theoretical part: Theory and Hypotheses
\begin{itemize}
\item[-] Briefly: Theoretical argument and argumentative structure
\item[-] Generate one hypothesis
\end{itemize}
\item Empirical part
\begin{itemize}
\item[3.1] Research Design
\begin{itemize}
\item[-] Discussion of sample and time span/case selection
\item[-] Discussion operationalization
\end{itemize}
\item[3.2] Discussion
\begin{itemize}
\item[-] Testing hypotheses using data visualization
\item[-] Describe and interpret your results
\end{itemize}
\end{itemize}

\item Conclusion
\begin{itemize}
\item[-] Summary and critical discussion
\item[-] Answering the research question
\end{itemize}	
\end{enumerate}


\end{document} 